\section{\texorpdfstring{Dynamical interpretation of the\\
local and collective memory contrast}{Dynamical interpretation of the local
and collective memory contrast}}
\label{sec:interpretation}

Section~\ref{sec:performance} establishes the controlled computational ordering
independently of the interpretation below. Here we ask which dynamical picture
is consistent with it after the tested initial states have been forgotten, not
which microscopic ingredient is uniquely causal.

\subsection{Working physical picture}

The simplest distinction concerns how the reservoir reaches its environment.
Under uniform local relaxation, every spin has its own direct relaxation
channel. Under collective relaxation, the spins share one channel. The same
assigned coupling weight is therefore spread over several separate paths in
one case and concentrated into a common path in the other. Some combinations
of spin excitations then reach the collective channel only after the
Hamiltonian has mixed them with the directly coupled combination. They are not
permanently protected: coherent dynamics continue to move information between
these combinations. Appendix~\ref{app:shared} gives the corresponding
mathematical formulation and states its precise limits. The same interference
principle appears in collective emission and in network states that a shared
environment cannot address directly~\cite{Dicke1954,Cabot2018}.

This gives the working picture. A shared environmental channel directly
couples fewer combinations. Hamiltonian mixing provides indirect relaxation
routes for the others, so the resulting slower readout-visible responses can
preserve a longer accessible input history.

\subsection{Scalar controls for the memory contrast}

Fixing \(\mathcal B\) controls assigned operator weight, but it does not
equalize realized jump activity, the relaxation spectrum, or the operating
point preferred by each environmental-coupling design. We therefore ask whether any one of
these scalar differences is sufficient to explain the principal memory
contrast. These controls are not intended to factorize the generator into
uniquely causal structural ingredients. Figure~\ref{fig:collective-case}(a) summarizes the
fixed-\(\mathcal B\), activity-matched, gap-matched, and independently selected
comparisons. The STM contrast is positive in all four protocols.

The first control asks whether the shared channel wins simply because it
produces a different amount of realized jump activity. Local and collective
rates are calibrated separately to the same time-averaged expected activity
on an input stream without task labels, then frozen before task scoring. Every
calibration passes and the held-out activity difference is unresolved.
Collective relaxation still gains \(3.306\) STM units, with a \(95\%\)
interval of \([1.393,5.220]\) and \(8/8\) paired wins. Different mean jump
activity cannot by itself explain the memory contrast.
Appendix~\ref{app:activity-control} gives the calibration protocol.

The second control asks whether the shared channel remembers longer simply
because it relaxes more slowly. Before
any task stream is generated, each local rate is therefore calibrated to the
collective generator's Liouvillian gap at the constant input \(s=0.5\). This
gap is the slowest nonzero decay rate of that constant-input evolution and
provides one measure of relaxation speed. All 24 matches fall within the
prescribed tolerance.

Gap matching reduces the STM and NARMA-10 contrasts by \(62.1\%\) and
\(48.2\%\), respectively, showing that this timescale explains an important
part of the effect. It does not remove the ordering. Collective relaxation
retains \(1.424\) STM units with an interval of \([0.800,2.049]\) and \(20/24\)
wins, and the NARMA-10 difference remains resolved.
Figure~\ref{fig:scalar-controls}(b) in Appendix~\ref{app:experiment1}
resolves the same comparison by input delay. Matching the local rate restores
much of its memory tail, but
collective relaxation remains higher at long delays. The bands are pointwise
\(95\%\) intervals across the 24 paired reservoirs.
Appendix~\ref{app:gap-control} gives the calibration protocol.

Finally, a common setting could accidentally favor one jump family. Local and
collective relaxation therefore choose their own drive, input duration,
dissipative multiplier, and ridge coefficient, with disjoint data for
selection and final testing. Collective relaxation still leads by \(3.945\)
STM units, with an interval of \([3.523,4.368]\) and \(24/24\) wins. The
contrast does not require a shared operating point. The search is bounded
rather than global. Appendix~\ref{app:selection-control} gives the selection
protocol and bounded-search scope.

These controls rule out unequal mean activity, one representative decay rate,
and a common operating point as sole explanations of the principal ordering
between local and collective relaxation. They do not equalize the complete spectrum,
input-dependent propagator products, non-normal transients, energy flow,
entropy production, implementation effort, or hardware cost, and they were
not repeated for the full cross-family task map.

\subsection{Positive dynamical support}

\RelaxationTheoryPriorWork

Because STM is obtained by linear regression on the measured feature
trajectory, a dynamical component can contribute only if the input excites it
and it remains visible in the readout span. Fixed-input mode overlap tests
readout visibility. The switched-response kernel additionally includes input
excitation. The slower collective modes remain Pauli-visible, and after the
representative timescale is matched the switched response is concentrated into
fewer observable patterns.

Two additional checks connect this response to the coupling-organization
interpretation. Across six
process designs, slower midpoint relaxation usually accompanies more STM. A
separate sweep mixes local and collective loss continuously. Its selected
interior point improves on local relaxation in every fresh pair, demonstrating
robustness to residual local damping along this path. Together, the diagnostics
support one coherent account of the endpoint contrast.
Appendix~\ref{app:interpretation} reports their uncertainty.

\subsection{Scope of the explanation}

Along the fixed-weight interpolation, every interior one-body lowering block
retains Kossakowski rank \(N\) while STM changes, so rank alone is insufficient.
Learned profiles and the interpolation change several coordinates together.
Two matched rank-one interventions isolate Kossakowski orientation. The
primary \(N=5\) study finds more STM along the equal-phase end of a
prespecified path and against five zero-overlap controls. In a separately
generated \(N=6\) replication, rotating the equal-phase vector to the real,
sign-balanced direction \((1,1,1,-1,-1,-1)\), orthogonal to the uniform
site profile of the input drive, lowers STM by \(2.271\) units
\([1.826,2.717]\), with
\(24/24\) paired wins for equal phase and exact two-sided sign-test
\(p=1.19\times10^{-7}\). Kossakowski rank, nonzero spectrum and trace,
assigned operator weight, sitewise diagonal weights, coefficient magnitudes,
and every within-pair processor and data choice remain fixed
(Figure~\ref{fig:phase-direction} in Appendix~\ref{app:phase-direction}).
These invariants are
therefore insufficient to determine memory: Kossakowski orientation relative
to the fixed driven processor is an operative coupling-organization
coordinate. This
does not distinguish alignment with the input from alignment with the
coherent dynamics, or make equal phase or a response diagnostic universal.

Dephasing follows a different contraction mechanism. In the connected primary
model, it removes the measured signal after washout, which accounts for its
near-zero STM. Appendix~\ref{app:interpretation} gives the driven-dephasing
argument. It should not be folded into the shared-channel explanation.
